% ============================================================================
%  Chapter 5 -- Number Representation: Integers                  (Lecture 05)
% ============================================================================
\chapter{Number Representation: Integers}
\label{ch:integers}

\epigraphline{I may count for little, it's true, said the One to the Zero, but
what are you worth? Nothing: nothing at all.}{Trilussa, \emph{Nummeri}, 1944}

\begin{chapgoals}
By the end of this chapter you will be able to
\begin{itemize}
\item convert integers and fractions between bases 2, 8, 10 and 16;
\item say which natural numbers fit in $n$ bits and recognise an overflow;
\item encode and decode signed integers in sign and magnitude and in two's complement;
\item add, subtract and negate in two's complement, and detect overflow;
\item explain why $0-1$ shows as \hex{FFFFFFFF} in a \riscv register;
\item choose the right \riscv instruction for signed and for unsigned data.
\end{itemize}
\end{chapgoals}

% ----------------------------------------------------------------------------
\section{Warm-up on the machine}

A marker stands on the screen, at position 2 of a horizontal line. The
program reads the joystick and, while \key{LEFT} is pressed, subtracts one
from the position at every frame (Listing~\ref{lst:int:warm}); the function
\code{draw} prints the position in hexadecimal with
\code{prg32\_console\_hex32}.

\asmfile[caption={Hold \key{LEFT} and watch the number (adapted from the laboratory
on state and input of the \prg textbook).},label=lst:int:warm]{code/l05/c03.S}

\begin{prgtry}
Hold \key{LEFT}. The value shown goes \texttt{00000002}, \texttt{00000001},
\texttt{00000000} \ldots\ and then \texttt{FFFFFFFF}, and the marker vanishes.
(On the board and in QEMU the console is drawn on the screen; \prgqt collects
it in its log.) Is it a bug? Is \hex{FFFFFFFF} the number 4\,294\,967\,295
or the number $-1$?
\end{prgtry}

The answer is: both. It depends on how the bits are read, and this chapter
gives the rules.

% ----------------------------------------------------------------------------
\section{Bases and conversions}

\subsection{Positional systems}

\begin{definition}{Positional representation}
Let $b\ge 2$ be an integer, the \term{base}. A string of digits
$d_{n-1}\,d_{n-2}\cdots d_1\,d_0$, with $0\le d_i\le b-1$ for every $i$,
represents in base $b$ the natural number
\begin{equation}\label{eq:int:positional}
N \;=\; \sum_{i=0}^{n-1} d_i\, b^{\,i}
  \;=\; d_{n-1}b^{\,n-1}+\dots+d_1 b+d_0 .
\end{equation}
The digit $d_{n-1}$ is the \term{most significant}, $d_0$ the \term{least
significant}.
\end{definition}

Four bases are used in computing (Table~\ref{tab:int:bases}). Every natural
number has exactly one representation in each base, apart from zeros added on
the left.

\begin{table}[htb]
\centering\small
\caption{The four bases of computing.}\label{tab:int:bases}
\begin{tabular}{@{}clll@{}}
\toprule
\textbf{Base $b$} & \textbf{Digits} & \textbf{Name} & \textbf{Thirteen is written}\\
\midrule
2 & 0, 1 & binary & \texttt{1101}\\
8 & 0, 1, \ldots, 7 & octal & \texttt{15}\\
10 & 0, 1, \ldots, 9 & decimal & \texttt{13}\\
16 & 0, \ldots, 9, A, \ldots, F & hexadecimal & \texttt{D}\\
\bottomrule
\end{tabular}
\end{table}

\subsection{Bases that are powers of two share the same bits}

If each digit is written in pure binary, the bits of a number are the same in
every base $b=2^{k}$: it is enough to group them from the right, $k$ at a
time, padding with zeros on the left. For the number 127:
\[
\texttt{01111111}_2 \;=\; \underbrace{\texttt{0111}}_{7}\,\underbrace{\texttt{1111}}_{F}
 \;=\; \hex{7F}
\qquad\text{and}\qquad
\texttt{01111111}_2 \;=\; \underbrace{\texttt{01}}_{1}\,\underbrace{\texttt{111}}_{7}\,\underbrace{\texttt{111}}_{7}
 \;=\; 177_8 .
\]
The reason is that in (\ref{eq:int:positional}) with $b=2$ the terms can be
collected $k$ at a time, and each group is a digit multiplied by a power of
$2^{k}$.

\subsection{Converting an integer: repeated division}

To write a natural number $N$ in base $r$: divide $N$ by $r$; the remainder is
the next digit, starting from the least significant; repeat on the quotient
until it is zero. The remainders are then read from the last to the first
(Table~\ref{tab:int:division}).

\begin{table}[htb]
\centering\small
\caption{The number 127 converted by repeated division.}\label{tab:int:division}
\begin{tabular}{@{}rrr@{}}
\multicolumn{3}{c}{$r=8$}\\\toprule
$N$ & $N/8$ & $N \bmod 8$\\\midrule
127 & 15 & 7\\ 15 & 1 & 7\\ 1 & 0 & 1\\\bottomrule
\multicolumn{3}{c}{$127 = 177_8$}
\end{tabular}\qquad
\begin{tabular}{@{}rrr@{}}
\multicolumn{3}{c}{$r=16$}\\\toprule
$N$ & $N/16$ & $N \bmod 16$\\\midrule
127 & 7 & 15 = F\\ 7 & 0 & 7\\\bottomrule
\multicolumn{3}{c}{$127 = \hex{7F}$}
\end{tabular}\qquad
\begin{tabular}{@{}rrr@{}}
\multicolumn{3}{c}{$r=2$}\\\toprule
$N$ & $N/2$ & $N \bmod 2$\\\midrule
127 & 63 & 1\\ 63 & 31 & 1\\ 31 & 15 & 1\\ 15 & 7 & 1\\ 7 & 3 & 1\\ 3 & 1 & 1\\ 1 & 0 & 1\\\bottomrule
\multicolumn{3}{c}{$127 = 1111111_2$}
\end{tabular}
\end{table}

The method is correct because, by (\ref{eq:int:positional}),
$N = r\,(d_{n-1}r^{\,n-2}+\dots+d_1) + d_0$: the remainder of the division of
$N$ by $r$ is $d_0$, and the quotient is the number whose digits are
$d_{n-1}\cdots d_1$, to which the same argument applies. In C the two
operations are \code{/} and \code{\%}; in \riscv, with the M extension, they
are \ins{divu} and \ins{remu}, and Chapter~\ref{ch:verm1p4} uses exactly this
algorithm to print a score.

\subsection{Converting a fraction: repeated multiplication}

A number $0\le N<1$ is written in base $r$ as
$N = d_{-1}r^{-1}+d_{-2}r^{-2}+\cdots$. Multiplying by $r$ gives
$d_{-1} + (d_{-2}r^{-1}+\cdots)$: the integer part of the product is the next
digit, starting from the most significant, and the procedure is repeated on
the fractional part (Table~\ref{tab:int:fraction}).

\begin{table}[htb]
\centering\small
\caption{Two fractions converted to base 2.}\label{tab:int:fraction}
\begin{tabular}{@{}lll@{}}
\multicolumn{3}{c}{0.875: exact}\\\toprule
Multiply & Fraction & Digit\\\midrule
$0.875\times 2$ & 0.75 & 1\\ $0.75\times 2$ & 0.5 & 1\\ $0.5\times 2$ & 0 & 1\\\bottomrule
\multicolumn{3}{c}{$0.875 = 0.111_2$}
\end{tabular}\qquad\qquad
\begin{tabular}{@{}lll@{}}
\multicolumn{3}{c}{0.8: periodic}\\\toprule
Multiply & Fraction & Digit\\\midrule
$0.8\times 2$ & 0.6 & 1\\ $0.6\times 2$ & 0.2 & 1\\ $0.2\times 2$ & 0.4 & 0\\ $0.4\times 2$ & 0.8 & 0\\ $0.8\times 2$ & 0.6 & 1 \ldots\\\bottomrule
\multicolumn{3}{c}{$0.8 = 0.\overline{1100}_2$}
\end{tabular}
\end{table}

\begin{prgwarn}
A fraction that is finite in base 10 can be periodic in base 2. With a finite
number of bits it must be cut: this is where the approximation errors of real
numbers come from (Chapter~\ref{ch:reals}).
\end{prgwarn}

% ----------------------------------------------------------------------------
\section{Natural numbers and overflow}

\subsection{Encoding strategies}

\begin{definition}{Representation}
A \term{representation} is a law $X = R(x)$ that associates a number $x$ with
a string of bits $X$. Here $x$ is the number that is meant; $X$ is the
pattern stored in the register.
\end{definition}

The number of bits depends on how many numbers are to be represented; it is
usually a multiple of 8, typically 16, 32 or 64. The law $R$ is chosen
according to the kind of number, natural, signed or rational, and a good law
is one that makes arithmetic operations and comparisons simple to build as
circuits.

\subsection{Natural numbers on $n$ bits}

With $n$ bits there are $2^{n}$ patterns, and the natural choice is to
associate each natural number from 0 to $2^{n}-1$ with the string of bits that
represents it in the binary positional system. With $n=4$ the interval is
$[0,15]$. The sizes handled by \riscv are in Table~\ref{tab:int:naturals}.

\begin{table}[htb]
\centering\small
\caption{Natural numbers on 8, 16 and 32 bits.}\label{tab:int:naturals}
\begin{tabular}{@{}rlll@{}}
\toprule
\textbf{Bits} & \textbf{Range} & \textbf{In \riscv} & \textbf{Maximum}\\
\midrule
8 & $0\ldots 2^{8}-1$ & byte (\ins{lbu}) & 255\\
16 & $0\ldots 2^{16}-1$ & halfword (\ins{lhu}) & 65\,535\\
32 & $0\ldots 2^{32}-1$ & word, register & 4\,294\,967\,295\\
\bottomrule
\end{tabular}
\end{table}

The number 4\,294\,967\,295 is the reading of \hex{FFFFFFFF} as a natural
number: the first half of the answer to the warm-up.

\subsection{Properties of binary strings}

Three properties follow at once from (\ref{eq:int:positional}).

\begin{enumerate}
\item \emph{Extension on the left with zeros does not change the value}:
  $5 = 101_2 = 00101_2$.
\item \emph{Multiplying by $2^{k}$ is a shift to the left by $k$ positions}:
  $00101_2\times 2 = 01010_2$, five becomes ten. Each digit $d_i$ moves to
  position $i+k$, where it weighs $2^{k}$ times as much.
\item \emph{Integer division by $2^{k}$ is a shift to the right by $k$
  positions}: $00101_2/2 = 00010_2$, five becomes two, and the remainder is
  lost.
\end{enumerate}

Shifts are much cheaper than multiplications, and compilers use them whenever
they can (Listing~\ref{lst:int:shifts}). The instruction
\code{slli t1, t1, 2} of Chapter~\ref{ch:information} multiplied an index by
four in this way. Shifts are also the basis of fixed-point arithmetic
(Chapter~\ref{ch:reals}).

\asmfile[caption={Shifts as multiplications and divisions.},label=lst:int:shifts]{code/l05/c04.S}

\subsection{Overflow}

\begin{definition}{Overflow}
An \term{overflow} occurs when the result of an operation is not contained in
the finite subset of numbers that the available bits can represent.
\end{definition}

Registers have a finite size, and there is a finite number of them. The sets
of numbers are infinite, and only a finite subset of each can be represented.
The arithmetic operators are closed on the whole set, not on the finite
subset. With 7 bits the interval $[0,127]$ is encoded: 100 and 38 are
representable, their sum 138 is not.

\subsection{Arithmetic on four bits}

Binary arithmetic follows the same column-by-column rules as decimal
arithmetic, with the single table $0+0=0$, $0+1=1$, $1+1=0$ with a carry
of~1.

\begin{center}\ttfamily
\begin{tabular}{@{}r@{\quad}r@{}}
\multicolumn{2}{c}{\textsf{\small Sum}}\\
6 + & 0110 +\\ 8 = & 1000 =\\\hline 14 & 1110
\end{tabular}\qquad
\begin{tabular}{@{}r@{\quad}r@{}}
\multicolumn{2}{c}{\textsf{\small Sum with overflow}}\\
14 + & 1110 +\\ 3 = & 0011 =\\\hline 17 & \textcolor{prgred}{1}\,0001
\end{tabular}\qquad
\begin{tabular}{@{}r@{\quad}r@{}}
\multicolumn{2}{c}{\textsf{\small Difference}}\\
11 - & 1011 -\\ 5 = & 0101 =\\\hline 6 & 0110
\end{tabular}\qquad
\begin{tabular}{@{}r@{\quad}r@{}}
\multicolumn{2}{c}{\textsf{\small Product}}\\
5 x & 0101 x\\ 3 = & 0011 =\\\hline 15 & 000\,1111
\end{tabular}
\end{center}

In the second sum a fifth bit is needed: the four-bit register keeps
\texttt{0001} and the carry is lost. Processors signal a lost carry in
different ways: many have a \emph{carry flag}; \riscv has none and lets the
program compute it with a comparison, as Section~\ref{sec:int:back} shows. A
product of two numbers of $n$ bits needs up to $2n$ bits: on RV32 the
instruction \ins{mul} gives the low 32 bits and \ins{mulhu} the high 32.

% ----------------------------------------------------------------------------
\section{Signed integers}

There are three ways of writing a minus sign with bits. Two are presented
here and the third, with a variant, in Chapter~\ref{ch:reals}.

\subsection{Sign and magnitude}

The most significant bit is the sign (1 for negative) and the other $n-1$ bits
are the magnitude:
\[
X = |x| + 2^{\,n-1}s,\qquad s=\begin{cases}1 & x<0\\ 0 & x\ge 0.\end{cases}
\]
The range goes from $-(2^{n-1}-1)$ to $+(2^{n-1}-1)$. Zero has two
representations, $+0$ and $-0$, so that only $2^{n}-1$ numbers are
represented. And arithmetic must work separately on sign and magnitude: to
add two numbers of opposite sign one must find the larger magnitude, subtract
the smaller from it, and give the result the sign of the operand with the
larger magnitude. The representation is simple for people and awkward for
circuits. It survives in the mantissa of floating-point numbers
(Chapter~\ref{ch:reals}).

\subsection{Two's complement}

\begin{definition}{Two's complement}
On $n$ bits, the \term{two's complement} representation of an integer $x$ with
$-2^{\,n-1}\le x\le 2^{\,n-1}-1$ is the natural number
\[
X=\begin{cases} x & \text{if } 0\le x<2^{\,n-1},\\[2pt]
2^{\,n}-|x| & \text{if } -2^{\,n-1}\le x<0,\end{cases}
\]
written in binary on $n$ bits. Equivalently, $X = x \bmod 2^{\,n}$.
\end{definition}

The range goes from $-2^{n-1}$ to $+2^{n-1}-1$, and zero has a single
representation. The most significant bit is 1 if and only if the number is
negative, because the negative numbers occupy the patterns from $2^{n-1}$ to
$2^{n}-1$. This is the representation used by every modern processor, \riscv
included: a 32-bit register read as a signed integer covers
$-2^{31}\ldots 2^{31}-1$, about minus two thousand million to plus two
thousand million. The range is asymmetric, with one more negative number than
positive ones, and this has a consequence that will appear shortly.

\subsection{Four bits, three readings}

Table~\ref{tab:int:readings} reads each of the sixteen patterns of four bits
in the three ways seen so far. When the most significant bit is 0 the three
readings agree; when it is 1 they differ.

\begin{table}[htb]
\centering\small
\caption{Four bits, three readings.}\label{tab:int:readings}
\begin{tabular}{@{}cccc@{\qquad\qquad}cccc@{}}
\toprule
\textbf{Bits} & \textbf{Unsigned} & \textbf{Sign-mag.} & \textbf{Two's c.} &
\textbf{Bits} & \textbf{Unsigned} & \textbf{Sign-mag.} & \textbf{Two's c.}\\
\midrule
\texttt{0000} & 0 & $+0$ & 0 & \texttt{1000} & 8 & $-0$ & $-8$\\
\texttt{0001} & 1 & $+1$ & $+1$ & \texttt{1001} & 9 & $-1$ & $-7$\\
\texttt{0010} & 2 & $+2$ & $+2$ & \texttt{1010} & 10 & $-2$ & $-6$\\
\texttt{0011} & 3 & $+3$ & $+3$ & \texttt{1011} & 11 & $-3$ & $-5$\\
\texttt{0100} & 4 & $+4$ & $+4$ & \texttt{1100} & 12 & $-4$ & $-4$\\
\texttt{0101} & 5 & $+5$ & $+5$ & \texttt{1101} & 13 & $-5$ & $-3$\\
\texttt{0110} & 6 & $+6$ & $+6$ & \texttt{1110} & 14 & $-6$ & $-2$\\
\texttt{0111} & 7 & $+7$ & $+7$ & \texttt{1111} & 15 & $-7$ & $-1$\\
\bottomrule
\end{tabular}
\end{table}

\begin{prgconcept}
The bits do not say which column is the right one. The program does, by
choosing its instructions.
\end{prgconcept}

\subsection{The number wheel}

With $n$ bits arithmetic is \emph{modulo $2^{n}$}: the number line closes into
a circle (Figure~\ref{fig:int:wheel}). Adding 1 moves one step clockwise,
whatever the reading. The two readings differ only in the place where the
wheel ``breaks''. In the unsigned reading it breaks between \texttt{1111} and
\texttt{0000}, and crossing that point is a carry: $15+1=0$. In the two's
complement reading it breaks between \texttt{0111} and \texttt{1000}, and
crossing that point is a signed overflow: $7+1=-8$. One step back from
\texttt{0000} is \texttt{1111}: that is why $0-1$ gives all ones.

\begin{figure}[htb]
\centering
\begin{tikzpicture}
  \draw[prgrule,line width=7mm] (0,0) circle (2.6);
  \foreach \k in {0,...,15}{
    \pgfmathsetmacro{\a}{90-\k*22.5}
    \pgfmathtruncatemacro{\s}{ifthenelse(\k>7,\k-16,\k)}
    \pgfmathtruncatemacro{\ba}{mod(int(\k/8),2)}
    \pgfmathtruncatemacro{\bb}{mod(int(\k/4),2)}
    \pgfmathtruncatemacro{\bc}{mod(int(\k/2),2)}
    \pgfmathtruncatemacro{\bd}{mod(\k,2)}
    \node[font=\ttfamily\scriptsize] at (\a:2.6) {\ba\bb\bc\bd};
    \node[font=\sffamily\small,text=prgsteel] at (\a:1.85) {\k};
    \ifnum\k>7
      \node[circle,fill=prgamber!30,inner sep=1pt,minimum size=7mm,font=\sffamily\small] at (\a:3.55) {$\s$};
    \else
      \node[circle,fill=prgaccent!20,inner sep=1pt,minimum size=7mm,font=\sffamily\small] at (\a:3.55) {$\s$};
    \fi
  }
  \draw[prgviolet,very thick,dashed] (101.25:2.1) -- (101.25:4.1) node[above left,lbl,text=prgviolet,align=right]{unsigned break:\\carry};
  \draw[prgred,very thick,dashed] (-78.75:2.1) -- (-78.75:4.1) node[below right,lbl,text=prgred,align=left]{signed break:\\overflow};
  \draw[->,thick,prgink] (80:1.1) arc (80:10:1.1) node[midway,below left,lbl]{$+1$};
\end{tikzpicture}
\caption{The number wheel for four bits. Inner numbers: unsigned value; ring:
bits; outer discs: two's complement value.}
\label{fig:int:wheel}
\end{figure}

The same adder serves both readings; only the place where the wheel breaks
changes.

\subsection{Two's complement arithmetic}

\begin{theorem}{Addition in two's complement}
Let $x$, $y$ and $x+y$ all lie in $[-2^{n-1},\,2^{n-1}-1]$, and let $X$, $Y$
be the representations of $x$ and $y$. Then the representation of $x+y$ is
$(X+Y)\bmod 2^{n}$: the sum of the two patterns as natural numbers, with the
carry out of the most significant bit discarded.
\end{theorem}

\noindent\emph{Proof.} $X\equiv x$ and $Y\equiv y \pmod{2^{n}}$ by definition,
hence $X+Y\equiv x+y \pmod{2^{n}}$; and the representation of $x+y$ is, by
definition, $(x+y)\bmod 2^{n}$. \hfill$\square$

\smallskip
Addition therefore needs no knowledge of the signs: the adder for natural
numbers already gives the right answer.

\begin{center}\ttfamily
\begin{tabular}{@{}r@{\quad}r@{}}
\multicolumn{2}{c}{\textsf{\small Opposite signs}}\\
2 + & 0010 +\\ -6 = & 1010 =\\\hline -4 & 1100
\end{tabular}\qquad\qquad
\begin{tabular}{@{}r@{\quad}r@{}}
\multicolumn{2}{c}{\textsf{\small Two negatives}}\\
-2 + & 1110 +\\ -3 = & 1101 =\\\hline -5 & \textcolor{prgcomment}{1}\,1011
\end{tabular}\qquad\qquad
\begin{tabular}{@{}r@{\quad}r@{}}
\multicolumn{2}{c}{\textsf{\small Product}}\\
2 x & 0010 x\\ -3 = & 1101 =\\\hline -6 & \textcolor{prgcomment}{001}\,1010
\end{tabular}
\end{center}

The same argument holds for multiplication: $XY\equiv xy\pmod{2^{n}}$, so the
low $n$ bits of a product are correct whatever the signs. In the third
example $2\times 13 = 26 = 1\,1010_2$, and the low four bits \texttt{1010}
read $-6$. This is why \riscv has a single \ins{mul} for signed and unsigned
operands.

\subsection{Negating a number}

By definition $-x$ is represented by $2^{n}-X$. Two methods compute it
without a subtraction.

\begin{description}[style=nextline,leftmargin=1.5em]
\item[Method 1: complement every bit, then add 1.]
  $6=\texttt{0110}$, complement \texttt{1001}, plus one: $\texttt{1010}=-6$.
  It works because the complement of every bit of $X$ is the number
  $(2^{n}-1)-X$, a subtraction from all ones in which no borrow can occur;
  adding 1 gives $2^{n}-X$.
\item[Method 2: scan from the right.] Starting from the right, copy the bits
  up to and including the first 1, and complement all the others.
  $5=\texttt{0101}$ gives $\texttt{1011}=-5$; $1=\texttt{0001}$ gives
  $\texttt{1111}=-1$. It is method~1 with the addition carried out mentally:
  the final $+1$ propagates through the trailing ones of the complement, which
  are the trailing zeros of the original.
\end{description}

Negating twice returns the original number, with one exception: the most
negative number, $-2^{n-1}$, has no positive partner, and its negation
overflows and returns the number itself.

\asmfile[caption={Negation in \riscv.},label=lst:int:neg]{code/l05/c05.S}

\begin{prgconcept}[Subtraction is addition]
$y - x = y + (2^{n}-x) \pmod{2^{n}}$. An adder and a complementer are the
only components needed to build addition, subtraction, multiplication and
division.
\end{prgconcept}

This is the engineering reason for two's complement: one circuit, the adder of
Chapter~\ref{ch:combinational}, does all the work. In the \riscv datapath
\ins{sub} is \ins{add} with the second operand complemented and a carry-in
of~1.

\subsection{Detecting overflow}

\begin{center}\ttfamily
\begin{tabular}{@{}r@{\quad}r@{}}
\multicolumn{2}{c}{\textsf{\small Unsigned, 8 bits}}\\
128 + & 1000 0000 +\\ 130 = & 1000 0010 =\\\hline 258 & \textcolor{prgred}{1} 0000 0010
\end{tabular}\qquad\qquad\qquad
\begin{tabular}{@{}r@{\quad}r@{}}
\multicolumn{2}{c}{\textsf{\small Signed, 8 bits}}\\
127 + & 0111 1111 +\\ 127 = & 0111 1111 =\\\hline 254 & \textcolor{prgred}{1}111 1110
\end{tabular}
\end{center}

In the first sum there is a carry out, and the register holds 2, the sum
modulo 256. In the second there is no carry out, but the result reads $-2$.

\begin{theorem}{Overflow conditions}
In an addition on $n$ bits:
\begin{enumerate}[nosep]
\item an \emph{unsigned} overflow occurs if and only if there is a carry out of
  the most significant bit;
\item a \emph{signed} overflow occurs if and only if two operands with the same
  sign give a sum with the opposite sign: two positives give a negative, or
  two negatives give a positive.
\end{enumerate}
The sum of two numbers of opposite sign never overflows.
\end{theorem}

The last statement holds because such a sum lies between the two operands,
both of which are representable. The two conditions are independent: the
same addition can trigger one, both or neither. In terms of circuits the
signed condition is equivalent to this: the carry \emph{into} the most
significant bit differs from the carry \emph{out} of it.

% ----------------------------------------------------------------------------
\section{Back to the machine}\label{sec:int:back}

\subsection{The warm-up, explained}

The variable \code{marker\_x} held 0. The instruction
\code{addi t2, t2, -1} moved one step back on a wheel with $2^{32}$ positions,
and the register now holds thirty-two ones: \hex{FFFFFFFF}. Read as unsigned
it is 4\,294\,967\,295; read in two's complement it is $-1$. The adder did the
same thing in both cases. The function \code{prg32\_console\_hex32} prints the
pattern, not an interpretation: it is neither right nor wrong. And the marker
disappears because \code{prg32\_gfx\_rect} is asked to draw at an enormous
$x$, or at $-1$: off the screen either way (Figure~\ref{fig:int:marker}). It
is not a bug of the processor; it is a missing decision in the program.

\begin{figure}[htb]
\centering
\includegraphics[width=0.5\textwidth]{lab2_step2_move.png}
\caption{The laboratory program: seven lamps for the bits of the buttons and a
marker.}
\label{fig:int:marker}
\end{figure}

\subsection{Signed or unsigned? The instruction chooses}

\riscv registers have no type. The same 32 bits are signed or unsigned only
in the eye of the instruction that reads them (Table~\ref{tab:int:signed}).

\begin{table}[htb]
\centering\small
\caption{Signed and unsigned instructions.}\label{tab:int:signed}
\begin{tabularx}{\textwidth}{@{}lllL@{}}
\toprule
\textbf{Operation} & \textbf{Signed} & \textbf{Unsigned} & \textbf{Why they differ}\\
\midrule
add, subtract & \ins{add}, \ins{addi}, \ins{sub} & the same & two's complement: one adder for both\\
multiply (low word) & \ins{mul} & the same & the low 32 bits do not depend on the sign\\
compare & \ins{slt}, \ins{slti} & \ins{sltu}, \ins{sltiu} & is \hex{FFFFFFFF} less than 0? yes if signed\\
branch & \ins{blt}, \ins{bge} & \ins{bltu}, \ins{bgeu} & the same question, as a jump\\
shift right & \ins{sra}, \ins{srai} & \ins{srl}, \ins{srli} & copy the sign bit, or insert zeros\\
load a byte & \ins{lb} & \ins{lbu} & sign extension or zero extension\\
divide, remainder & \ins{div}, \ins{rem} & \ins{divu}, \ins{remu} & magnitudes are read differently\\
\bottomrule
\end{tabularx}
\end{table}

\subsection{Keeping the marker on the screen}

The decision that was missing is taken in Listing~\ref{lst:int:clamp}: after
moving, the position is \term{clamped} to the interval $[0, 304]$.

\asmfile[caption={Clamping the position (laboratory on state and input, step 3).},label=lst:int:clamp]{code/l05/c08.S}

\ins{bgez} reads \reg{t2} as a signed number: it tests the sign bit.
\ins{ble} is a signed comparison too: $-1$ is less than 304, whereas with
\ins{bleu} the pattern \hex{FFFFFFFF} would be far greater. The constant
\code{-SPEED} is an immediate: twelve bits in two's complement, from $-2048$
to $+2047$.

\begin{prgconcept}
A game does not detect overflow after the fact: it keeps its values inside a
range where overflow cannot happen.
\end{prgconcept}

\subsection{Detecting overflow without flags}

\riscv has no carry flag and no overflow flag: fewer hidden states, a simpler
processor. When a program does need to know, it asks with a comparison
(Listing~\ref{lst:int:ovf}). For unsigned numbers, a sum that has wrapped
round is always smaller than either operand. For signed numbers, adding a
negative number must make the sum smaller, and adding a non-negative one must
not. On the wheel, \ins{sltu} detects the crossing of the top and the signed
test the crossing of the bottom.

\asmfile[caption={The two overflow tests suggested by the \riscv specification.},label=lst:int:ovf]{code/l05/c09.S}

\subsection{Sign extension: the same number on more bits}

\begin{theorem}{Sign extension}
A number in two's complement on $n$ bits is written on $m>n$ bits by copying
its sign bit into all the new positions on the left.
\end{theorem}

\noindent\emph{Proof.} If $x\ge 0$ the sign bit is 0 and the statement is the
extension with zeros of natural numbers. If $x<0$ its pattern on $n$ bits is
$X=2^{n}-|x|$ and the pattern wanted is $2^{m}-|x| = X + (2^{m}-2^{n})$; the
number $2^{m}-2^{n}$ is exactly $m-n$ ones followed by $n$ zeros.
\hfill$\square$

\begin{center}\ttfamily
\begin{tabular}{@{}r@{\quad}l@{\quad}l@{}}
\textsf{$-1$ on 4 bits} & \phantom{1111\,}1111 & = -1\\
\textsf{$-1$ on 8 bits} & \textcolor{prgred}{1111}\,1111 & = -1\\
\textsf{$+5$ on 4 bits} & \phantom{0000\,}0101 & = +5\\
\textsf{$+5$ on 8 bits} & \textcolor{prgsteel}{0000}\,0101 & = +5\\
\end{tabular}
\end{center}

Extending with zeros would turn \texttt{1111} ($-1$) into \texttt{0000\,1111}
($+15$). The instruction \code{addi t2, t2, -1} stores $-1$ in twelve bits,
and the processor sign-extends it to thirty-two before adding. \ins{lb}
sign-extends the byte that it loads, \ins{lbu} fills with zeros: the choice
depends on what the byte means.

% ----------------------------------------------------------------------------
% Generated by tools/textbook/export_listings.js from decks/plan.js and code/l05.S: do not edit.
\section{The program of this chapter}\label{sec:lab:l05}

The listings of this chapter are parts of one program, \fname{code/l05.S} in the
repository of the course, printed in full in Listing~\ref{lst:lab:l05}. The
automatic tests of the course material build it and run it on the headless
runner of \prgqt{}, with scripted input, and check the frame;
Figure~\ref{fig:lab:l05} is the frame that the test wrote.

\begin{description}[style=nextline,leftmargin=1.5em]
\item[On the screen.] The 32 bits of marker\_x as lamps. It starts at 2 (one lamp); each frame with LEFT held subtracts 1: after 2, 1 and 0 every lamp lights up, 0xFFFFFFFF, which is -1.
\item[What to press.] Hold LEFT for a few frames.
\end{description}

\begin{figure}[htb]
\centering
\includegraphics[width=0.5\textwidth]{lab_l05.png}
\caption{\fname{l05.S} on \prgqt{}.}\label{fig:lab:l05}
\end{figure}

\asmfile[caption={The complete program of the chapter, \fname{code/l05.S}.},label=lst:lab:l05]{code/l05/program.S}

To run it on \prgqt{} the commands are those of Listing~\ref{lst:lab:l05:run};
Section~\ref{sec:prg:testing} gives the same steps for QEMU and for the board.

\begin{shell}[caption={Building \fname{l05.S} and running it on \prgqt{}.},label=lst:lab:l05:run]
# 1. start PRG32-QT on this computer (it listens on port 8080)

# 2. in the PRG32 folder, with ESP-IDF active: build the cartridge
python3 -m prg32 cartridge build ~/arc12/code/l05.S --portable \
    --entry-prefix l05 --name l05 --out l05.prg32

# 3. send it to PRG32-QT and start it
python3 -m prg32 esp32c6 upload-and-run l05.prg32 \
    --url http://localhost:8080 --slot cart0

# 4. look, press, compare with the previous slide
\end{shell}


% ----------------------------------------------------------------------------
\section{Check your understanding}

\begin{quiz}{Which number does the 8-bit pattern \texttt{1111 1100} represent in two's complement?}
\item 252
\item $-4$
\item $-3$
\item $-124$
\end{quiz}

\begin{quiz}{On 8 bits, which addition produces a signed overflow?}
\item $100 + 27$
\item $100 + 28$
\item $-100 + 28$
\item $-1 + 1$
\end{quiz}

\begin{quiz}{\reg{t0} holds \hex{FFFFFFFF} and \reg{t1} holds 1. After
\code{slt t2, t0, t1} and \code{sltu t3, t0, t1}:}
\item \reg{t2} = 1, \reg{t3} = 1
\item \reg{t2} = 0, \reg{t3} = 0
\item \reg{t2} = 1, \reg{t3} = 0
\item \reg{t2} = 0, \reg{t3} = 1
\end{quiz}

\begin{quiz}{How is 0.625 written in binary?}
\item 0.101
\item 0.110
\item 0.625
\item 0.011
\end{quiz}

\section{Exercises}

\begin{exercise}[paper first, then \prg]
\begin{enumerate}[nosep]
\item Convert 2026 to binary, octal and hexadecimal; convert 0.1 to binary up
  to eight fractional bits.
\item Write $-100$ on 8 bits in sign and magnitude and in two's complement.
\item Compute on 8 bits in two's complement: $90+60$, $-90-60$, $90-60$.
  Which overflow?
\item In the clamp code replace \ins{bgez} and \ins{ble} with their unsigned
  versions, predict the behaviour, then run it.
\item Write an \code{update} that adds \hex{10000000} to a counter at each
  frame and lights a red rectangle when the signed sum overflows.
\end{enumerate}
\expected{On paper: $2026 = \hex{7EA} = 111\,1110\,1010_2$. On \prg: the red
rectangle flashes once every sixteen frames, when the counter passes from
\hex{7\ldots} to \hex{8\ldots}.}
\end{exercise}

\section{Glossary}

\begin{glossarytable}
base (radix) & base\\
natural number & numero naturale\\
signed integer & numero relativo (intero con segno)\\
most significant bit & bit pi\`u significativo\\
remainder & resto\\
shift & scorrimento (shift)\\
two's complement & complemento a due (alla base)\\
sign and magnitude & segno e modulo\\
overflow & overflow (traboccamento)\\
carry / borrow & riporto / prestito\\
sign extension & estensione del segno\\
to clamp & limitare a un intervallo\\
\end{glossarytable}

\section{Conclusions and reading}

Every base has its advantages: base 2 for circuits, base 16 for people who
read bits. $n$ bits represent $2^{n}$ numbers: which ones depends on the
encoding law. Two's complement turns subtraction into addition, so that one
adder does everything. Overflow is the price of finite registers: it is
detected, or designed away. In \riscv the bits have no type: the instruction
chooses between signed and unsigned.

\paragraph{Reading.} \citet{patterson2020}, sections 2.4 and 3.1--3.2;
\emph{Playing with RISC-V}, chapter~6, ``State and Input''
\citep{prg32textbook}; \citet{stallings2021}, the chapters on number systems
and on computer arithmetic.
